\documentclass[12pt]{article}
\usepackage[margin=1in]{geometry}
\usepackage{booktabs}
\usepackage{amsmath}
\usepackage{hyperref}
\usepackage{setspace}
\onehalfspacing

\title{Problem Set 7 \\ \large Econ 5253 --- Spring 2026}
\author{Yeyang Zhou}
\date{March 31, 2026}

\begin{document}
\maketitle

\section*{Question 5: Dropping Missing Observations}

After loading \texttt{wages.csv} (2,246 observations), we drop all rows
where either \texttt{hgc} or \texttt{tenure} is missing (17 observations
total), leaving \textbf{2,229 observations} for analysis.

\section*{Question 6: Summary Statistics and Missing Data Discussion}

Table~\ref{tab:summary} presents summary statistics for the cleaned sample.

\begin{table}[ht]
\centering
\caption{Summary Statistics ($N = 2{,}229$)}
\label{tab:summary}
\begin{tabular}{lrrrrr}
\toprule
 & Count & Mean & Std.\ Dev. & Min & Max \\
\midrule
logwage & 1669   & 1.6252  & 0.3855 & 0.0049 & 2.2615 \\
hgc     & 2229   & 13.1014 & 2.5243 & 0.0000 & 18.0000 \\
college & 2229   & 0.2378  & 0.4258 & 0 & 1 \\
tenure  & 2229   & 5.9706  & 5.5072 & 0.0000 & 25.9167 \\
age     & 2229   & 39.1516 & 3.0620 & 34.0000 & 46.0000 \\
married & 2229   & 0.6420  & 0.4795 & 0 & 1 \\
\bottomrule
\end{tabular}
\end{table}

\noindent\textbf{Missing rate for \texttt{logwage}:}
Of the 2,229 observations, 560 have missing log wages, for a missing
rate of \textbf{25.1\%}.

\medskip
\noindent\textbf{MCAR, MAR, or MNAR?}
The \texttt{logwage} variable is most likely \textbf{Missing Not at
Random (MNAR)}. Women with lower wages---e.g.\ those in informal or
part-time employment---are less likely to have their wages recorded.
This means the probability of missingness is correlated with the
unobserved value of \texttt{logwage} itself, which rules out both
MCAR (missingness independent of all data) and MAR (missingness
dependent only on observed covariates).

\section*{Question 7: Imputation Methods and Regression Results}

We estimate the following model under four imputation strategies:
\begin{equation}
\log w_i = \beta_0 + \beta_1\,\text{hgc}_i + \beta_2\,\text{college}_i
         + \beta_3\,\text{tenure}_i + \beta_4\,\text{tenure}_i^2
         + \beta_5\,\text{age}_i + \beta_6\,\text{married}_i + \varepsilon_i
\end{equation}

The four methods are: (1) \textbf{complete cases} (listwise deletion,
assumes MCAR); (2) \textbf{mean imputation} (replace missing
\texttt{logwage} with the sample mean); (3) \textbf{predicted value
imputation} (impute using fitted values from the complete-cases
regression, consistent with MAR); and (4) \textbf{multiple imputation}
(5 imputed datasets with added residual noise, pooled via Rubin's rules).

Table~\ref{tab:regression} presents all four sets of estimates.
The true value of $\beta_1 = 0.093$.

\begin{table}[ht]
\centering
\caption{Returns to Schooling under Alternative Imputation Methods}
\label{tab:regression}
\begin{tabular}{lllll}
\hline
               & Complete   & Mean       & Predicted  & Multiple    \\
               & Cases      & Imputation & Imputation & Imputation  \\
\hline
Intercept      & 0.6567***  & 0.8490***  & 0.6567***  & 0.7089***   \\
               & (0.1301)   & (0.1029)   & (0.0991)   & (0.1133)    \\
hgc ($\hat{\beta}_1$) & 0.0624*** & 0.0497*** & 0.0624*** & 0.0590*** \\
               & (0.0054)   & (0.0043)   & (0.0042)   & (0.0048)    \\
college        & $-$0.1452*** & $-$0.1682*** & $-$0.1452*** & $-$0.1250*** \\
               & (0.0345)   & (0.0257)   & (0.0247)   & (0.0283)    \\
tenure         & 0.0495***  & 0.0382***  & 0.0495***  & 0.0498***   \\
               & (0.0052)   & (0.0039)   & (0.0037)   & (0.0043)    \\
tenure$^2$     & $-$0.0016*** & $-$0.0013*** & $-$0.0016*** & $-$0.0016*** \\
               & (0.0003)   & (0.0002)   & (0.0002)   & (0.0002)    \\
age            & 0.0004     & 0.0002     & 0.0004     & 0.0002      \\
               & (0.0027)   & (0.0022)   & (0.0021)   & (0.0024)    \\
married        & 0.0220     & 0.0268**   & 0.0220*    & 0.0143      \\
               & (0.0177)   & (0.0136)   & (0.0131)   & (0.0150)    \\
$R^2$          & 0.2084     & 0.1473     & 0.2767     & 0.2158      \\
$N$            & 1669       & 2229       & 2229       & 2229        \\
\hline
\multicolumn{5}{l}{\small Standard errors in parentheses.\
  $^*p<0.10$,\ $^{**}p<0.05$,\ $^{***}p<0.01$}
\end{tabular}
\end{table}

\noindent\textbf{Discussion of $\hat{\beta}_1$ across models:}

All four estimates of $\hat{\beta}_1$ fall below the true value of 0.093,
consistent with MNAR: lower-wage observations are missing
disproportionately, compressing the observed variation in log wages and
attenuating the estimated returns to schooling.

\textbf{Complete cases (0.0624):} Dropping missing observations yields a
$\hat{\beta}_1$ notably below the true value. Because low-wage workers are
more likely to be missing, the MNAR pattern pulls the estimate downward.

\textbf{Mean imputation (0.0497):} The lowest estimate among the four
methods. Replacing every missing value with the same constant artificially
compresses the variance of the dependent variable, attenuating all slope
coefficients toward zero. Mean imputation is generally a poor strategy
and should be avoided.

\textbf{Predicted value imputation (0.0624):} Identical to the
complete-cases estimate, as expected---the imputed values are exactly the
fitted values from that regression, so OLS recovers the same coefficients.
The standard errors are artificially small because residual variance is
understated.

\textbf{Multiple imputation (0.0590):} The closest estimate to the true
value among the four methods. By adding stochastic noise to imputed
values, MI restores variance in the dependent variable and produces more
credible standard errors via Rubin's rules. Despite the MNAR nature of
the data, MI performs best in practice and is the preferred approach.

Overall, the results illustrate the classic attenuation bias from poorly
handling missing data. Mean imputation is clearly the worst strategy;
multiple imputation is the most reliable.

\section*{Question 8: Project Progress}

For my final project, I am using \textbf{Compustat firm financials data},
which contains annual accounting and financial information for publicly
traded U.S.\ firms, including variables such as total assets, revenues,
capital expenditures, leverage, and profitability measures across a long
panel of firms and years.

My primary modeling approach will be \textbf{Ordinary Least Squares
(OLS)} regression to examine the relationship between key firm financial
characteristics and outcomes of interest such as investment or firm
growth. Depending on the structure of the data, I may also explore panel
fixed effects specifications to control for time-invariant firm
heterogeneity.

\end{document}
